% Macros comuns aos cards de divulgação.
% =============================================================================
% marinha-comum.tex — identidade visual compartilhada pela apresentação e pela
% apostila da capacitação. Segue o Manual de Identidade Visual da Marinha do
% Brasil (dez/2025): cores institucionais (p. 23) e tipografia (pp. 24-30).
%
% Compilar sempre com XeLaTeX (fontspec). Os caminhos são relativos à pasta
% do documento (apresentacao/ ou apostila/).
% =============================================================================

% ---------- Caminhos ---------------------------------------------------------
\newcommand{\raizcapacitacao}{..}
\newcommand{\imagens}{\raizcapacitacao/assets/imagens}
\newcommand{\fontes}{\raizcapacitacao/assets/fontes}

% ---------- Cores institucionais (Manual, p. 23) -----------------------------
\usepackage{xcolor}
\usepackage{graphicx}
\definecolor{azulMB}{HTML}{050F41}   % Azul Marinha  · Pantone 2768C
\definecolor{ouroMB}{HTML}{FAB932}   % Amarelo Ouro  · Pantone 1235C
\definecolor{verdeMB}{HTML}{079551}  % Verde Brasil  · Pantone 355C
% Tons de apoio derivados das institucionais (fundos de caixas e diagramas)
\colorlet{azulClaro}{azulMB!8}
\colorlet{azulMedio}{azulMB!55}
\colorlet{ouroClaro}{ouroMB!18}
\colorlet{verdeClaro}{verdeMB!12}
\definecolor{cinzaTexto}{HTML}{3A3F55}
\definecolor{cinzaSuave}{HTML}{F3F4F8}
\definecolor{vermelhoAlerta}{HTML}{B3261E}
\colorlet{vermelhoClaro}{vermelhoAlerta!10}

% ---------- Tipografia (Manual, pp. 24-30) -----------------------------------
% Principal: Octin College (títulos) e alternativa Bebas Neue Pro são fontes
% comerciais. Usamos a Bebas Neue livre (OFL), da mesma família de desenho,
% nos títulos, e a Montserrat (tipografia de apoio oficial, OFL) no texto.
\usepackage{fontspec}
\setmainfont{Montserrat}[
  Path=\fontes/montserrat/, Extension=.otf,
  UprightFont=*-Regular, BoldFont=*-SemiBold,
  ItalicFont=*-Italic, BoldItalicFont=*-SemiBoldItalic]
\setsansfont{Montserrat}[
  Path=\fontes/montserrat/, Extension=.otf,
  UprightFont=*-Regular, BoldFont=*-SemiBold,
  ItalicFont=*-Italic, BoldItalicFont=*-SemiBoldItalic]
\setmonofont{FiraMono}[
  Path=\fontes/fira-mono/, Extension=.otf,
  UprightFont=*-Regular, BoldFont=*-Bold, Scale=0.9]
\newfontfamily\titulofonte{BebasNeue-Regular}[
  Path=\fontes/bebas-neue/, Extension=.ttf, LetterSpace=2.0]
\newfontfamily\montblack{Montserrat-Black}[Path=\fontes/montserrat/, Extension=.otf]
\newfontfamily\montlight{Montserrat-Light}[Path=\fontes/montserrat/, Extension=.otf]

% ---------- Pacotes comuns ----------------------------------------------------
\usepackage{tikz}
\usetikzlibrary{positioning,arrows.meta,calc,shapes.geometric,fit,backgrounds,shadows.blur,decorations.pathmorphing}
\usepackage{booktabs,tabularx,array}
\usepackage{fontawesome5}
\usepackage{listings}
\usepackage{csquotes}

% ---------- Código (listings) -------------------------------------------------
\lstdefinestyle{terminal}{
  basicstyle=\ttfamily\small\color{cinzaTexto},
  columns=fixed, basewidth=0.6em, keepspaces=true,
  breaklines=true, breakatwhitespace=true,
  showstringspaces=false, upquote=true,
  commentstyle=\color{verdeMB},
  keywordstyle=\color{azulMB}\bfseries,
  morekeywords={git,glab,claude},
  morecomment=[l]{\#},
  literate=
    {á}{{\'a}}1 {à}{{\`a}}1 {ã}{{\~a}}1 {â}{{\^a}}1
    {é}{{\'e}}1 {ê}{{\^e}}1 {í}{{\'i}}1 {ó}{{\'o}}1
    {õ}{{\~o}}1 {ô}{{\^o}}1 {ú}{{\'u}}1 {ç}{{\c c}}1
    {Á}{{\'A}}1 {É}{{\'E}}1 {Í}{{\'I}}1 {Ó}{{\'O}}1
    {Ú}{{\'U}}1 {Ç}{{\c C}}1 {Ã}{{\~A}}1 {Õ}{{\~O}}1
    {→}{{$\rightarrow$}}1 {—}{{---}}1 {·}{{$\cdot$}}1
}
\lstdefinestyle{saida}{style=terminal, keywordstyle=, commentstyle=,
  basicstyle=\ttfamily\small\color{cinzaTexto}}

% ---------- Diagramas de histórico Git (TikZ) ---------------------------------
\tikzset{
  commit/.style={circle, draw=#1, fill=white, line width=1.6pt,
    minimum size=7mm, inner sep=0pt, font=\small\bfseries\color{#1}},
  commit/.default=azulMB,
  commitm/.style={circle, draw=ouroMB!80!black, fill=ouroMB, line width=1.6pt,
    minimum size=8mm, inner sep=0pt, font=\small\bfseries\color{azulMB}},
  hist/.style={line width=1.8pt, color=#1},
  hist/.default=azulMB,
  rotulo/.style={rounded corners=2pt, fill=#1, text=white,
    font=\scriptsize\bfseries, inner sep=3pt},
  rotulo/.default=azulMB,
  caixa/.style={rounded corners=3pt, draw=#1, line width=1pt, fill=#1!8,
    align=center, font=\small, inner sep=5pt, text=azulMB},
  caixa/.default=azulMB,
  seta/.style={-{Stealth[length=3mm]}, line width=1.2pt, color=azulMB},
  setafina/.style={-{Stealth[length=2mm]}, line width=0.8pt, color=azulMedio},
}

% ---------- Marca do projeto (texto) -------------------------------------------
% Usada onde a versão interna tinha o logotipo da Marinha, que não entra na
% versão pública. \marca[cor das letras]{tamanho em pt}
\newcommand{\marca}[2][azulMB]{{\titulofonte\fontsize{#2}{#2}\selectfont
  \textcolor{#1}{GIT}\textcolor{ouroMB}{\,·\,}\textcolor{#1}{IA}\textcolor{ouroMB}{\,·\,}\textcolor{#1}{DOCS}}}

\usepackage{polyglossia}
\setdefaultlanguage[variant=brazilian]{portuguese}
\pagestyle{empty}
\usetikzlibrary{fadings}
\hyphenpenalty=10000
\exhyphenpenalty=10000
\newcommand{\seta}{\raisebox{0.35ex}{\fontsize{22}{22}\selectfont\faLongArrowAltRight}}
% \numero{coordenada}{número}{rótulo}
\newcommand{\numero}[3]{%
  \node[anchor=north west, text=ouroMB, font=\titulofonte\fontsize{40}{40}\selectfont] at (#1) {#2};
  \node[anchor=north west, text=white, font=\small, text width=3.9cm] at ($(#1)+(0.05,-1.45)$) {#3};}
% Mini grafo: dois históricos sem relação unidos por um merge
\newcommand{\grafo}[1][1]{%
  \begin{scope}[scale=#1, every node/.style={transform shape}]
    \foreach \l/\x in {A/0,B/1.3,C/2.6,D/3.9}{\node[commit=verdeMB, fill=azulMB, font=\small\bfseries\color{white}] (\l) at (\x,0.9) {\l};}
    \draw[hist=verdeMB] (A)--(B)--(C)--(D);
    \foreach \l/\x in {X/0,Y/1.9}{\node[commit=white, fill=azulMB, font=\small\bfseries\color{white}] (\l) at (\x,-0.9) {\l};}
    \draw[hist=white] (X)--(Y);
    \node[commitm] (M) at (5.6,0) {M};
    \draw[hist=verdeMB] (D) to[out=0,in=135] (M);
    \draw[hist=white] (Y) to[out=0,in=-135] (M);
  \end{scope}}
